\section{An alphabet-uniform embedding by substring histograms}
\label{upper:section}

We now construct an embedding on the whole set of words of bounded length.
For a finite alphabet $\Sigma$, let $\Sigma^{\le d}$ include all words of
length at most $d$, including the empty word.  Write $\operatorname{ED}(x,y)$
for the minimum number of single-symbol insertions, deletions, and
substitutions taking $x$ to $y$, each at unit cost.  Intermediate words may
have arbitrary length.  All logarithms in this section are natural.

\begin{theorem}[Uniform upper bound]
\label{upper:uniform}\label{thm:uniform-upper}
There are absolute constants $C>0$ and $d_0\ge3$ such that, for every finite
alphabet $\Sigma$ with $|\Sigma|\ge2$ and every integer $d\ge d_0$, there
are an integer $M\ge1$ and a deterministic map
$F:\Sigma^{\le d}\longrightarrow\ell_1^M$ satisfying
\[
 \operatorname{ED}(x,y)\le \|F(x)-F(y)\|_1
 \le \exp\!\bigl(C\sqrt{\log d\,\log\log d}\bigr)
       \operatorname{ED}(x,y)
 \qquad(x,y\in\Sigma^{\le d}).
\]
Here $\ell_1^M$ is real, and both $C$ and $d_0$ are independent of $\Sigma$.
\end{theorem}

The construction develops the block-and-shift recursion of Ostrovsky and
Rabani \cite{OstrovskyRabani2007}.  They compare multisets of overlapping
shifted substrings using embeddings of shorter strings, and their
introduction also identifies extensions to larger alphabets and varying
lengths.  We give the estimates for these extensions explicitly.

The recursive step represents a block by lists of overlapping windows.
After embedding the shorter windows, we record how many of them fall in
each part of a random partition. Averaging over every partition outcome
gives one deterministic $\ell_1$ map. A pair of distant blocks has a
scale at which every cross-pair of windows is distant, so their histograms
are often disjoint. For expansion, a deletion followed by an insertion
shifts most blocks by one position, allowing most windows to be paired
identically. These observations prove the theorem first at power lengths.
We then construct a map using the actual window lengths and analyze its
expansion both through an arbitrary longest common subsequence and through
moves. A final variant uses total variation in the power-length recursion.

\subsection{Subsequence deficits and local moves}

A common subsequence is an increasing matching of equal symbols; its
maximum size is $\operatorname{LCS}(x,y)$.  For arbitrary words put
\[
 q(x,y)=\max(|x|,|y|)-\operatorname{LCS}(x,y),
\]
and for words of the same length $n$ put
$r_n(x,y)=n-\operatorname{LCS}(x,y)$.  Thus $r_n$ counts the unmatched
positions on either one side of a longest common subsequence.

\begin{lemma}[Metric comparisons]\label{upper:metric}
Let $\Delta(x,y)$ be the minimum number of insertions and deletions when
substitutions are forbidden.  For all finite words,
\[
 \Delta(x,y)=|x|+|y|-2\operatorname{LCS}(x,y),\qquad
 \operatorname{ED}(x,y)\le\Delta(x,y)\le2\operatorname{ED}(x,y),
\]
and
\begin{equation}\label{upper:unequal-comparison}
 q(x,y)\le\operatorname{ED}(x,y)\le2q(x,y).
\end{equation}
In particular, for $|x|=|y|=n$,
\begin{equation}\label{upper:deficit-comparison}
 r_n(x,y)=\tfrac12\Delta(x,y),\qquad
 r_n(x,y)\le\operatorname{ED}(x,y)\le2r_n(x,y).
\end{equation}
\end{lemma}
\begin{proof}
In a script of $k$ edits, at least $|x|-k$ original symbols survive
unchanged and in order.  Applying this also to the reverse script gives
$q(x,y)\le k$.  Deleting the letters outside a longest common subsequence
and inserting the missing target letters uses
$|x|+|y|-2\operatorname{LCS}(x,y)\le2q(x,y)$ edits.  Conversely, the
survivors of any insertion--deletion script form a common subsequence,
so that script has at least this many edits.  This proves the formula for
$\Delta$.  Ordinary edits allow insertions and deletions, while replacing
each substitution by one deletion and one insertion proves the factor-two
comparison.  These alignment identities are the insertion--deletion
specialization of the trace formulation of Wagner and Fischer
\cite[Sections~3 and~5]{WagnerFischer1974}.
\end{proof}

\begin{lemma}[Moves and their restrictions]\label{upper:move}
Two words of length $n\ge0$ can be joined by exactly $r_n(x,y)$ moves, each
deleting one symbol and inserting one symbol while preserving length $n$.
For a single move, restrictions to any common interval of $a$ positions
have deficit at most one, and hence edit distance at most two.

For a fixed partition of the positions into nonempty consecutive blocks, at most
two blocks contain a deletion or insertion endpoint.  In any other block
of length $m$, the two subwords agree, or the prefix of length $m-1$ of
one agrees with the suffix of length $m-1$ of the other.
\end{lemma}
\begin{proof}
Protect the occurrences of a longest common subsequence in $x$, labeling
them by their positions in $y$.  Delete an unprotected symbol and insert
a missing target symbol in its prescribed order among the protected
ones, then protect it.  After $r_n(x,y)$ repetitions every target position
is protected, and the word is $y$.

For one move, let $p$ be its deleted position in $x$ and $q$ its inserted
position in $y$.  If $p<q$, positions outside $[p,q]$ match identically,
while positions $p+1,\ldots,q$ of $x$ match positions $p,\ldots,q-1$ of
$y$.  Restricting this matching to the same interval on both sides loses
at most one position on each side, so its size is at least $a-1$.
Interchanging the words handles $p>q$; when $p=q$ only one position
changes.  The same description shows that a block avoiding both endpoints
is unchanged or shifted by one position, in the precise sense stated.
\end{proof}

\begin{lemma}[Common and fresh suffixes]\label{upper:padding}
For equal-length words $u,v$ of length $m$ and any common suffix $z$,
\begin{equation}\label{upper:common-suffix}
 r_{m+|z|}(uz,vz)=r_m(u,v).
\end{equation}
If $\ast$ occurs in neither $x$ nor $y$, and
$N\ge\max(|x|,|y|)$, then, writing
$\widehat x=x\ast^{N-|x|}$ and $\widehat y=y\ast^{N-|y|}$,
\begin{equation}\label{upper:fresh-padding}
 \operatorname{LCS}(\widehat x,\widehat y)
 =\operatorname{LCS}(x,y)+N-\max(|x|,|y|),\qquad
 r_N(\widehat x,\widehat y)=q(x,y).
\end{equation}
The second identity includes empty inputs.
\end{lemma}
\begin{proof}
An optimal subsequence matching of two words ending in the same symbol
can be chosen to match their last positions.  If an optimal matching uses
either last position, its last such pair can be replaced by the pair of
last positions; if it uses neither, that pair can be appended.  Removing
the final pair shows
$\operatorname{LCS}(ua,va)\le1+\operatorname{LCS}(u,v)$.
Appending the pair to an optimal matching of the prefixes proves the
reverse inequality.
Iteration over $z$ proves \eqref{upper:common-suffix}.

For fresh-symbol padding, matched original symbols form a common
subsequence of $x,y$, and matched $\ast$ symbols form a final part of
size at most $\min(N-|x|,N-|y|)$.  These two maxima are attained
consecutively.  Their sum gives \eqref{upper:fresh-padding}.
\end{proof}

The recursion at equal lengths will use a common suffix with arbitrary
symbols.  The final passage to varying lengths will use the fresh symbol
in \eqref{upper:fresh-padding}.

\subsection{A finite histogram component}

The next construction turns a distance between shorter words into a
distance between lists of such words.  We use the standard cut
representation of finite $\ell_1$ metrics
\cite[Section~4]{LinialLondonRabinovich1995}. The proof below includes the
summability needed when the original target has infinitely many coordinates.

\begin{lemma}[Finite random partitions]\label{upper:partitions}
Let $X$ be finite and
$\rho(u,v)=\|f(u)-f(v)\|_1$ for a map $f:X\to\ell_1$ into real sequence
space.  For every $\alpha>0$ there is a probability distribution with
finite support on partitions $\pi$ of $X$ such that
\begin{equation}\label{upper:collision}
 \Pr[\pi(u)=\pi(v)]=e^{-\alpha\rho(u,v)}.
\end{equation}
Here $\pi(u)$ denotes the part containing $u$.  For a finite list $L$
in $X$, let $H_L^\pi$ be its count vector on the parts of $\pi$, with
repetitions counted.  Then
\begin{equation}\label{upper:histogram-map}
 \mathcal H_\alpha(L)=\bigoplus_\pi \Pr[\pi]\,H_L^\pi
\end{equation}
is a deterministic finite-dimensional map whose distance equals
$\mathbb E\|H_L^\pi-H_{L'}^\pi\|_1$.
\end{lemma}
\begin{proof}
For real $a,b$,
$|a-b|=\int_{\mathbb R}|\mathbf 1_{a>t}-\mathbf 1_{b>t}|\,dt$.
Apply this to every coordinate of $f$ and group the threshold cuts by
the subset of $X$ they select.  After discarding the empty and full cuts,
\[
 \rho(u,v)=\sum_{\varnothing\ne A\subsetneq X}
 a_A|\mathbf 1_A(u)-\mathbf 1_A(v)|,\qquad a_A\ge0.
\]
Every nontrivial cut separates some unordered pair.  Consequently
$\sum_Aa_A\le\sum_{\{u,v\}\subset X}\rho(u,v)<\infty$.
This also justifies the grouping when $f$ has infinitely many coordinates;
all sums and integrals are nonnegative, and there are only finitely many
subsets of $X$.

Select each cut $A$ independently with probability $1-e^{-\alpha a_A}$,
and partition $X$ by the membership bits in the selected cuts.  Points
$u,v$ share a part exactly when no separating cut is selected.  Taking
the product of these probabilities gives \eqref{upper:collision}.
There are finitely many cut selections.  Finally, additivity of the
$\ell_1$ norm on disjoint coordinates proves the distance identity for
\eqref{upper:histogram-map}.
\end{proof}

\begin{lemma}[Histogram estimates]\label{upper:histogram}
Let $s\ge1$ be an integer and let $L,L'$ be lists of size $s$ in the
set of Lemma~\ref{upper:partitions}, and set
$T=\mathbb E\|H_L^\pi-H_{L'}^\pi\|_1$.
If every cross-pair has $\rho$-distance at least $a$, then
\begin{equation}\label{upper:histogram-lower}
 T\ge2s(1-s^2e^{-\alpha a}).
\end{equation}
If a pairing of list occurrences leaves $t$ occurrences unpaired on
each side, then
\begin{equation}\label{upper:coupling}
 T\le2t+2\sum_{(u,v)\ {\rm paired}}(1-e^{-\alpha\rho(u,v)})
   \le2t+2\alpha\sum_{(u,v)\ {\rm paired}}\rho(u,v).
\end{equation}
Writing $P_L^\pi=H_L^\pi/s$ and
$\operatorname{TV}(P,Q)=\tfrac12\|P-Q\|_1$, the exact normalizations are
\begin{equation}\label{upper:normalizations}
 \|H_L^\pi-H_{L'}^\pi\|_1
 =s\|P_L^\pi-P_{L'}^\pi\|_1
 =2s\,\operatorname{TV}(P_L^\pi,P_{L'}^\pi).
\end{equation}
\end{lemma}
\begin{proof}
The probability of any cross-collision is at most $s^2e^{-\alpha a}$
by a union bound.  With no cross-collision the two count vectors have
disjoint supports and mass $s$ each, hence distance $2s$.  This proves
\eqref{upper:histogram-lower}.  For the upper bound, each unpaired
occurrence costs one in the triangle inequality.  A paired pair costs
two if its parts differ and zero otherwise.  Take expectations using
\eqref{upper:collision} and then use $1-e^{-z}\le z$ for $z\ge0$.
The normalization identities follow from the definitions.
\end{proof}

Thus a lower estimate for one component follows from separation of all
cross-pairs, while an upper estimate follows from a pairing of occurrences.
Both estimates concern the same map \eqref{upper:histogram-map}, whose
partition distribution is fixed on the full finite domain before any
pair of input words is considered.

\subsection{Power lengths and count histograms}\label{upper:powers}

\begin{proof}[Proof of Theorem~\ref{upper:uniform}]
Fix integers $b\ge2$ and $H\ge1$, put $n=b^H$, and fix a nonempty finite
alphabet $\Gamma$.  We construct maps $F_k$ at lengths $b^k$, $0\le k\le H$,
with
\begin{equation}\label{upper:power-normalization}
 r_{b^k}(x,y)\le\|F_k(x)-F_k(y)\|_1
             \le\Lambda_k r_{b^k}(x,y).
\end{equation}
At length one, $a\mapsto\tfrac12e_a$ has distance exactly $r_1$, so
$\Lambda_0=1$.

Index positions from zero and write $x_{[a,b)}$ for the subword in
positions $a,\ldots,b-1$.  Suppose $1\le k\le H$, set
$m=b^{k-1}$ and $N=bm$, and assume \eqref{upper:power-normalization}
at level $k-1$.  Split each length-$N$ word into $b$ consecutive blocks
of length $m$.  At each dyadic scale $s\le m$, block $j$ supplies the
list
\[
 W_{j,s}(x)=
 \bigl(x_{[jm+r,\,jm+r+m-s+1)}\bigr)_{r=0}^{s-1},
 \qquad 0\le j<b.
\]
These are its $s$ windows of length $m-s+1$, counted with multiplicity.
Choose one fixed suffix $z_s\in\Gamma^{s-1}$ and embed a window $u$
by $F_{k-1}(uz_s)$.  By \eqref{upper:common-suffix}, the resulting
distance $\rho_s$ satisfies
\[
 r_{m-s+1}(u,v)\le\rho_s(u,v)
       \le\Lambda_{k-1}r_{m-s+1}(u,v).
\]
Apply Lemma~\ref{upper:partitions} on the full domain
$\Gamma^{m-s+1}$ at rate
\begin{equation}\label{upper:lambda-local}
 \alpha_s=\frac{4\log(s+1)}s.
\end{equation}
Let $T_{j,s}(x,y)$ be the expected count-histogram distance of the two
lists, and let $T_{k,s}=\sum_jT_{j,s}$.  Define $F_k$ as four times the
direct sum of these histogram maps over $j$ and $s$.  Thus
\[
 \|F_k(x)-F_k(y)\|_1=4\sum_sT_{k,s}(x,y).
\]
Every map in this construction is finite-dimensional, since each level
takes a finite direct sum of the maps in Lemma~\ref{upper:partitions}.

\paragraph{Separation.}
Suppose corresponding blocks have deficit $q>0$.  The largest dyadic
integer $s\le(q+1)/2$ is available and satisfies
\begin{equation}\label{upper:power-scale}
 \frac{q+1}{4}<s\le\frac{q+1}{2},\qquad s\ge\frac q4.
\end{equation}
Every cross-pair $u,v$ of their windows satisfies
\[
 r_{m-s+1}(u,v)
 \ge(m-s+1)-\operatorname{LCS}(\text{the two full blocks})
 =q-s+1\ge s.
\]
The union bound in \eqref{upper:histogram-lower} is therefore at most
$s^2/(s+1)^4\le1/2$, and this component is at least $s\ge q/4$.
Common subsequences chosen separately in the blocks concatenate, so the
sum of their deficits is at least $r_N(x,y)$.  It follows that
$\|F_k(x)-F_k(y)\|_1\ge r_N(x,y)$.

\paragraph{Expansion along a move.}
In a nonexceptional block from Lemma~\ref{upper:move}, the two window
lists have at least $s-1$ identical occurrences.  For a shifted block,
pair starts displaced by one; for an unchanged block, pair equal starts.
Equation~\eqref{upper:coupling} bounds its contribution by $2$.
In an exceptional block, pair windows at the same starts.  Their
deficits are at most one by Lemma~\ref{upper:move}, including after the
common suffix is appended, so one such block costs at most
$2s\alpha_s\Lambda_{k-1}=8\log(s+1)\Lambda_{k-1}$.
There are at most two exceptional blocks.  At every scale,
\begin{equation}\label{upper:power-move}
 T_{k,s}(x,y)\le2b+16\log(s+1)\Lambda_{k-1}
 \quad\text{for one move.}
\end{equation}
There are at most $1+\lfloor\log_2m\rfloor$ scales.  Summing this estimate,
applying the factor $4$, and following the $r_N(x,y)$-move path gives
\begin{equation}\label{upper:power-recurrence}
 \Lambda_k\le C_1\log^2(n+2)(b+\Lambda_{k-1}),\qquad
 \Lambda_H\le b\bigl(C_2\log^2(n+2)\bigr)^H
\end{equation}
with absolute constants.  Indeed, after enlarging
$A=C_1\log^2(n+2)$ so that $A\ge2$, induction gives
$\Lambda_k\le b(2A)^k$.

\paragraph{Parameters and varying lengths.}
For sufficiently large $d$, take
\begin{equation}\label{upper:parameters}
 H=\left\lfloor\sqrt{\frac{\log d}{\log\log d}}\right\rfloor,\qquad
 b=\left\lceil d^{1/H}\right\rceil,\qquad n=b^H.
\end{equation}
Here $d\le n$ and, since $b\le2d^{1/H}$,
$\log n\le\log d+H\log2=O(\log d)$.  Also
$H$ has order $\sqrt{\log d/\log\log d}$, so
\[
 \log\Lambda_H
 \le\log b+H\log\!\bigl(C_2\log^2(n+2)\bigr)
 =O\!\bigl(\sqrt{\log d\,\log\log d}\bigr).
\]
All constants are absolute.  Work over
$\Gamma=\Sigma\cup\{\ast\}$ with $\ast\notin\Sigma$ and append $\ast$
to every input until its length is $n$.  By
\eqref{upper:unequal-comparison} and \eqref{upper:fresh-padding},
$F(x)=2F_H(\widehat x)$ satisfies
\[
 \operatorname{ED}(x,y)\le\|F(x)-F(y)\|_1
 \le2\Lambda_H\operatorname{ED}(x,y).
\]
Absorbing the factor two into the absolute exponential constant proves
the theorem, including the empty word.
\end{proof}

\subsection{Actual window lengths and two expansion estimates}
\label{upper:variants}\label{upper:aligned}

The power-length construction padded each window to the preceding power.
We now use maps at the actual window lengths. This gives one recursion
whose expansion can be proved in two ways: directly from an arbitrary
longest common subsequence, or from the move estimate already used above.
Both arguments analyze the same deterministic histogram map.

Fix an integer cap $d\ge2$, an integer $B\ge2$, and a nonempty finite
alphabet $\Gamma$. Construct maps $\Phi_n$ on $\Gamma^n$, $1\le n\le d$,
whose distances lie between $\operatorname{ED}$ and
$E_n\operatorname{ED}$, starting with $\Phi_1(a)=\tfrac12e_a$ and $E_1=1$.
For $n>1$, partition the positions into $b_n=\min(B,n)$ nonempty blocks
of nearly equal lengths, with boundaries
\[
 0=p_0<p_1<\cdots<p_{b_n}=n,\qquad
 \ell_j=p_{j+1}-p_j\le\lceil n/B\rceil.
\]
For each dyadic $s\le\ell_j$, use the list of $s$ windows
$x_{[p_j+r,\,p_j+r+\ell_j-s+1)}$, $0\le r<s$.
Embed each window by the map at its actual length. Writing
\[
 D'=\max_{1\le t\le\lceil n/B\rceil}E_t,
\]
its embedded distance satisfies
$\operatorname{ED}\le\rho\le D'\operatorname{ED}$.
Every child length is strictly less than $n$. Use count histograms at rate
\begin{equation}\label{upper:aligned-rate}
 \alpha_s=\lambda/s,\qquad \lambda=4\log(2d),
\end{equation}
and define $\Phi_n$ as six times their direct sum. Write $T_{j,s}$ for
the expected count distance. Equivalently, each component before the factor
six is $s$ times the expected $\ell_1$ distance of probability histograms,
by~\eqref{upper:normalizations}.

\paragraph{Separation.}
If corresponding blocks have edit distance $k_j>0$, then
$k_j\le\ell_j$ by substitution. The largest dyadic
$s\le(k_j+2)/3$ is available and satisfies $s\ge k_j/6$ and
$k_j-2(s-1)\ge s$. Deleting the positions omitted from the two windows
shows that every cross-window edit distance is at least
$k_j-2(s-1)$. The probability of any cross-collision is at most
$d^2e^{-\lambda}\le1/2$, so $T_{j,s}\ge s\ge k_j/6$.
Blockwise edit scripts concatenate; hence the sum of the block distances
dominates the whole-word distance. The factor six makes $\Phi_n$
noncontracting.

\paragraph{Corresponding segments in an arbitrary alignment.}
Fix an arbitrary longest common subsequence matching of $x,y$, of
size $n-k$, where $k=r_n(x,y)$.  Let $t_j$ be the number of matched
source positions strictly before $p_j$.  Choose target boundaries
\[
 0=q_0\le q_1\le\cdots\le q_{b_n}=n
\]
so that exactly $t_j$ matched target positions precede $q_j$.
For an internal boundary, take $q_j$ to be the position of target
match number $t_j+1$ if it exists, and take $q_j=n$ otherwise.
This gives nondecreasing boundaries and permits repeated boundaries
and empty target segments.  The matching pairs exactly the matched
positions of $[p_j,p_{j+1})$ with those of $[q_j,q_{j+1})$.

Let $A_j,B_j$ be the numbers of unmatched source and target positions
in the prefixes ending at these boundaries.  Since
$p_j=t_j+A_j$, $q_j=t_j+B_j$, and $0\le A_j,B_j\le k$,
\begin{equation}\label{upper:boundary-shift}
 |q_j-p_j|\le k.
\end{equation}
Let $a_j,b_j$ count unmatched positions in the source block and its
aligned target segment, and put $e_j=a_j+b_j$.  These segments
partition the two words, so
\begin{equation}\label{upper:unmatched-total}
 \sum_j e_j=2k.
\end{equation}
In a matched pair belonging to these segments, the relative indices
measured from $p_j,q_j$ differ by at most $e_j$: their numbers of
preceding matched positions agree, and their preceding unmatched
counts are bounded by $a_j,b_j$.

\paragraph{Pairing the original window lists.}
Put $L=\ell_j-s+1$.  Pair the source window starting at $p_j+r$
with the target window starting at $q_j+r$ precisely when both belong
to the original fixed-block lists.  Thus
\[
 0\le r<s,\qquad p_j\le q_j+r<p_j+s.
\]
The two intervals of allowed starts overlap in
$\max(0,s-|q_j-p_j|)$ positions.  Exactly
$\min(s,|q_j-p_j|)$ windows on each side remain unpaired.

Consider a paired window.  At most $a_j\le e_j$ of its source
positions are unmatched.  For any other source position, its
partner's relative index differs by at most $e_j$.  Therefore only
matched source positions within $e_j$ positions of an endpoint can
have partners outside the target window; there are at most $2e_j$
of them.  The windows have a common subsequence of length at least
$\max(0,L-3e_j)$, and \eqref{upper:deficit-comparison} gives
\begin{equation}\label{upper:aligned-window}
 \operatorname{ED}(\text{paired windows})\le6e_j.
\end{equation}
This reasoning also covers a target window extending outside its aligned
segment.  Every partner of a matched source-block position lies in that
segment, and the argument counts precisely those partners that remain
inside the target window.  It requires no condition on the other
target positions.

Using the partial pairing in \eqref{upper:coupling}, the shorter-map
expansion, and \eqref{upper:aligned-rate}, we obtain
\begin{equation}\label{upper:aligned-component}
 T_{j,s}(x,y)
 \le2|q_j-p_j|+2s\frac{\lambda}{s}D'(6e_j)
 =2|q_j-p_j|+12\lambda D'e_j.
\end{equation}
Each block has at most $L_d=1+\lfloor\log_2d\rfloor$ scales.
Summing~\eqref{upper:aligned-component}, using
\eqref{upper:boundary-shift} and~\eqref{upper:unmatched-total}, and
applying the factor six gives
\[
 \|\Phi_n(x)-\Phi_n(y)\|_1
 \le6L_d(2Bk+24\lambda D'k)
 \le C_6\log^2(d+1)(B+D')\operatorname{ED}(x,y).
\]
The last inequality uses $k=r_n(x,y)\le\operatorname{ED}(x,y)$.
The chosen alignment serves only to analyze two images; every partition
and coordinate of $\Phi_n$ was fixed on the entire domain beforehand.

\paragraph{An alternative expansion estimate from moves.}
For a single move, a nonexceptional block has $s-1$ identical windows
and therefore contributes at most two. In either of the at most two
exceptional blocks, windows at equal starts have edit distance at most
two. Their embedded distance is at most $2D'$, so
\eqref{upper:coupling} gives a block cost at most $4s\alpha_sD'=4\lambda D'$.
Thus, at each scale,
\begin{equation}\label{upper:ed-move}
 \sum_jT_{j,s}(x,y)\le2B+8\lambda D'
 =2B+32\log(2d)D'.
\end{equation}
Summing scales, multiplying by six, and following the
$r_n(x,y)\le\operatorname{ED}(x,y)$ moves gives the alternative expansion estimate
\begin{equation}\label{upper:ed-move-recurrence}
 E_n\le C_3\log(2d)\bigl(B+\log(2d)D'\bigr).
\end{equation}
The arbitrary-alignment argument established this order of expansion
without decomposing the comparison into moves.

\paragraph{Depth and varying lengths.}
Both estimates have the form $E_n\le A(B+D')$ with
$A=C\log^2(2d)$ for an absolute $C$. The identity
\begin{equation}\label{upper:ceiling-identity}
 \left\lceil\frac{\lceil n/B\rceil}{B}\right\rceil
 =\left\lceil\frac{n}{B^2}\right\rceil
\end{equation}
follows because both sides are the least integer $t$ with $n\le B^2t$.
It bounds the recursion depth by $h=\lceil\log_Bd\rceil$.
The induction used in~\eqref{upper:power-recurrence} gives expansion at most
\begin{equation}\label{upper:ceiling-recurrence}
 B\bigl(C_7\log^2(2d)\bigr)^h.
\end{equation}
For sufficiently large $d$, take
$B=\lceil\exp(\sqrt{\log d\,\log\log d})\rceil$. The logarithm
of this bound is then $O(\sqrt{\log d\,\log\log d})$.

Finally, work over $\Gamma=\Sigma\cup\{\ast\}$ and append the fresh symbol
to length $d$. Erasing $\ast$ from a padded edit script projects each
operation to at most one edit, and~\eqref{upper:deficit-comparison}
and~\eqref{upper:fresh-padding} give
\[
 \operatorname{ED}(x,y)\le\operatorname{ED}(\widehat x,\widehat y)
 \le2r_d(\widehat x,\widehat y)=2q(x,y)
 \le2\operatorname{ED}(x,y).
\]
Thus $\Phi_d(\widehat x)$ proves Theorem~\ref{upper:uniform} on all shorter
words, including the empty word, with an additional expansion factor two.

\subsection{Total-variation normalization}\label{upper:tv}

A final variant retains the power lengths $n_i=b^i$ and the common-suffix
window embeddings of Section~\ref{upper:powers}. Construct $V_i$ with
distances between $r_{n_i}$ and $D_i r_{n_i}$, starting with
$V_0(a)=\tfrac12e_a$ and $D_0=1$. At scale $s$, use rate
\begin{equation}\label{upper:tv-rate}
 \alpha_s=\frac{\log(2s^2)}s
\end{equation}
and $s$ times expected total variation of probability histograms.
This component is half the count distance. For block deficit $q>0$,
choose $s$ as in~\eqref{upper:power-scale}. Every cross-window deficit
is at least $s$, so the probability of any cross-collision is at most
$s^2e^{-\log(2s^2)}=1/2$. The component is at least $s/2>q/8$;
eight times the summed map is therefore noncontracting for $r_{n_i}$.

Along a move, a nonexceptional block costs at most one, while an
exceptional block costs at most $s\alpha_sD_{i-1}$. Consequently the
unscaled cost at each scale is at most
\begin{equation}\label{upper:tv-move}
 b+2\log(2s^2)D_{i-1},
\end{equation}
and the move path gives
\begin{equation}\label{upper:tv-recurrence}
 D_i\le C_5\log(2n_i)\bigl(b+\log(2n_i)D_{i-1}\bigr).
\end{equation}
Choose $b,H,n$ as in~\eqref{upper:parameters}. Bounding the logarithmic
factors by those at $n=b^H$ and using~\eqref{upper:power-recurrence}
yields $D_H\le b(C\log^2(2n))^H$ for an absolute $C$.
The parameter calculation in the power-length proof gives the required
exponent. Its fresh-symbol padding and final multiplication by two apply
to $V_H$, proving the same theorem on all shorter words.
